\chapter{Anonimato e privatezza su Internet}

Abbiamo chiuso il capitolo precedente con uno sniffer sulla porta 443: un nome leggibile, e tutto il resto Application Data. Ma a
chiunque sul percorso restano leggibili gli \textbf{indirizzi}, il \textbf{nome del server}, le \textbf{dimensioni} e i \textbf{tempi}.
Uno sniffer impara quindi chi parla con chi, e quando, anche quando non legge nulla di ciò che si dicono. TLS nasconde \emph{che cosa}
si dice; in questo capitolo vediamo se qualcosa può nascondere \emph{a chi} lo si dice. Il percorso: da chi ci si nasconde; che cosa
lascia trapelare ogni livello; gli strumenti; ciò che nessuno strumento può risolvere.

% ══════════════════════════════════════════════════════════════
\section{Anonimi rispetto a chi}
% ══════════════════════════════════════════════════════════════

Alice vuole leggere il sito di un attivista politico. Non vuole che:

\begin{table}[H]
\centering
\begin{tabular}{lcc}
\toprule
\textbf{Alice non vuole che} & \textbf{HTTP} & \textbf{HTTPS} \\
\midrule
il sito conosca il suo indirizzo IP & esposto & esposto \\
il suo ISP sappia quale sito visita & esposto & esposto \\
il suo ISP legga ciò che scambia & esposto & \textbf{nascosto} \\
\bottomrule
\end{tabular}
\caption{Tre cose che Alice vuole nascondere.}
\end{table}

Il suo indirizzo sorgente è in ogni datagramma che invia, e l'indirizzo del sito pure. TLS risolve una sola delle tre righe; alla fine
del capitolo la tabella avrà due colonne in più.

\begin{figure}[H]
\centering
\begin{tikzpicture}[every node/.style={font=\scriptsize\sffamily}]
  \node (a) at (0,0) {\icon[0.8cm]{laptop}}; \node[below=-1pt of a, font=\footnotesize\bfseries\sffamily] {Alice};
  \node[netcloud, minimum width=1.9cm, minimum height=1cm] (l) at (2.9,0) {rete locale};
  \node[netcloud, minimum width=1.6cm, minimum height=1cm] (i) at (5.9,0) {ISP};
  \node[netcloud, minimum width=2.2cm, minimum height=1cm] (t) at (9.1,0) {reti di transito};
  \node (s) at (12.3,0) {\icon[0.45cm]{server}}; \node[below=-1pt of s, font=\footnotesize\bfseries\sffamily] {il sito};
  \draw[wired] (a) -- (l); \draw[wired] (l) -- (i); \draw[wired] (i) -- (t); \draw[wired] (t) -- (s);
  \node[draw, fill=llink!50, rounded corners] (r) at (5.9,1.5) {resolver};
  \draw[wired, dashed] (i) -- (r);
  \node[anchor=west, text width=3.6cm] at (6.6,1.5) {ogni nome che si cerca, dal proprio indirizzo};
  \node[text width=1.8cm, align=center, text=netred] at (0,-1.25) {invia tutto};
  \node[text width=2.4cm, align=center, text=netred] at (2.9,-1.25) {MAC, indirizzo, ogni frame inviato};
  \node[text width=2.4cm, align=center, text=netred] at (5.9,-1.25) {indirizzo, ogni destinazione, ogni nome};
  \node[text width=2.4cm, align=center, text=netred] at (9.1,-1.25) {entrambi gli indirizzi, dimensioni e tempi};
  \node[text width=2.2cm, align=center, text=netred] at (12.3,-1.25) {il vostro indirizzo, tutto ciò che chiedete};
  \draw[netorange, thick, {Stealth}-{Stealth}] (0.3,-2.0) -- node[below, text=netorange] {un osservatore di entrambi gli estremi: confronta ciò che entra con ciò che esce} (12.1,-2.0);
\end{tikzpicture}
\caption{Ogni osservatore sta in un punto diverso e vede una parte diversa; il più pericoloso vede entrambi gli estremi.}
\end{figure}

Non esiste essere anonimi \emph{in generale}: solo anonimi \textbf{rispetto a qualcuno}.

\subsection{Leggere e bloccare sono minacce diverse}

\begin{itemize}
    \item Un \textbf{osservatore} vuole leggere o collegare: chi siete, con chi parlate, che cosa dite.
    \item Un \textbf{censore} non vuole né l'una né l'altra cosa: vuole fermare la conversazione, e gli basta riconoscerla.
          \begin{itemize}
            \item \textbf{Per indirizzo}: scarta ogni pacchetto verso un elenco di indirizzi IP (il filtro di pacchetti per cui abbiamo
                  scritto regole).
            \item \textbf{Per nome}: risponde alle query DNS con bugie, o resetta le connessioni la cui SNI è in un elenco.
            \item \textbf{Per forma}: l'ispezione profonda dei pacchetti riconosce un protocollo anche quando non ne legge un byte.
          \end{itemize}
\end{itemize}

Nascondere chi si è non ferma un censore; nascondere che aspetto ha il traffico potrebbe farlo.

\subsection{Quattro cose diverse}

\dfn{Quattro proprietà}{
    \begin{itemize}
      \item \textbf{Confidenzialità}: nessuno, a parte i due estremi, può leggere il contenuto.
      \item \textbf{Anonimato}: non si può essere identificati all'interno di un insieme di soggetti, l'\textbf{insieme di anonimato}
            (\emph{anonymity set}).
      \item \textbf{Pseudonimato}: si agisce sotto un identificatore che non è il proprio nome, ma che è \textbf{sempre lo stesso}.
      \item \textbf{Non collegabilità} (\emph{unlinkability}): nessuno può dire se due messaggi, o due visite, vengono dallo stesso
            soggetto.
    \end{itemize}
    Sono indipendenti: uno pseudonimo è collegabile per costruzione, è a questo che serve. La confidenzialità non dice nulla sulle altre
    tre, ed è l'unica che TLS fornisce.
}

\warning{Lo pseudonimo che un tribunale ha smascherato}{
    Un \emph{remailer} riceve la vostra posta, toglie il vostro indirizzo e la ripubblica sotto uno pseudonimo. Chi può scoprire che
    eravate voi? Chi possiede la sua tabella da pseudonimi a indirizzi, e chi può costringere il gestore a consegnarla: negli anni '90 un
    tribunale fece esattamente questo. Lo pseudonimato è anonimato solo finché la corrispondenza resta segreta (un remailer di
    ``tipo 1''). I remailer \emph{cypherpunk} non tenevano né tabelle né log: un server sequestrato non aveva nulla da dare. Gli utenti
    ne concatenavano tre, ogni messaggio cifrato per il successivo: incontreremo di nuovo questa forma.
}

\warning{Un lucchetto e un login}{
    Accedete a un sito in HTTPS. Quale delle quattro proprietà avete? Solo la confidenzialità, e solo rispetto al percorso. Il sito ha il
    vostro nome, e il vostro ISP sa quale sito è. TLS autentica il server; il vostro login autentica voi: insieme sono l'\textbf{opposto}
    dell'anonimato. Scambiare la cifratura per anonimato è l'errore più comune sulla privatezza in Internet.
}

% ══════════════════════════════════════════════════════════════
\section{Che cosa lascia trapelare ogni livello}
% ══════════════════════════════════════════════════════════════

\begin{table}[H]
\centering
\small
\begin{tabular}{*{5}{>{\raggedright\arraybackslash}p{2.95cm}}}
\toprule
 & \textbf{collegamento} & \textbf{IP} & \textbf{TCP} & \textbf{TLS} \\
\midrule
primo pacchetto & MAC sorg., MAC dest. & \texttt{192.168.1.23} $\rightarrow$\newline \texttt{143.225.161.30} & \texttt{51734} $\rightarrow$ \texttt{443} &
Client Hello, SNI = \texttt{www.unina.it} \\
record successivo & header di collegamento & header IP & header TCP & tipo e lunghezza del record; il resto cifrato \\
\midrule
chi lo legge & la LAN, l'access point & ogni router, il sito & NAT, firewall & chiunque sul percorso \\
\bottomrule
\end{tabular}
\caption{Una richiesta HTTPS, letta dall'esterno: richiesta, cookie e pagina li legge solo il server; dimensione e tempo li vede chiunque.}
\end{table}

TLS cifra il \textbf{proprio} payload. Ogni header sotto di esso viaggia esattamente come è stato scritto, e anche la parte cifrata dice
la propria lunghezza e il momento in cui è stata inviata.

\dfn{Un header è scritto per essere letto}{
    Sotto il proprio livello un payload è opaco; un header è l'opposto: esiste per essere letto lungo la strada. Lo switch deve leggere
    l'indirizzo MAC, ogni router l'indirizzo IP, il NAT le porte, il resolver il nome. Un livello può quindi cifrare solo ciò che gli sta
    \textbf{sopra}; il suo header resta in chiaro per i suoi dispositivi. TLS sta sopra TCP: dalle porte in giù tutto è in chiaro, e deve
    esserlo. Ogni livello lascia trapelare esattamente ciò che gli serve per consegnare; scendendo nella pila, la perdita cambia
    proprietario.
}

\subsection{Che cosa si dice all'applicazione}

\begin{itemize}
    \item In cima spesso si dice chi si è: un account, un login, un indirizzo di consegna, un numero di telefono. Nessun livello sotto può
          annullarlo: il miglior anonimato di rete finisce alla pagina di login.
    \item Senza login c'è il \textbf{cookie}: un numero che il server vi ha dato e che rimandate con ogni richiesta. Quel numero è uno
          pseudonimo: collega tutte le vostre visite senza bisogno del vostro nome. Come abbiamo visto con i cookie: date il vostro nome una
          volta, e l'intera cronologia ne acquista uno.
\end{itemize}

\begin{figure}[H]
\centering
\begin{tikzpicture}[every node/.style={font=\scriptsize\sffamily}]
  \foreach \y/\t in {1.3/un giornale, 0/un negozio, -1.3/una clinica} {
    \node[draw, fill=cloudbg, rounded corners, minimum width=2.2cm, minimum height=0.7cm] at (0,\y) {\t};
    \node[draw, fill=netorange!25, minimum width=2.4cm, minimum height=0.5cm, font=\scriptsize\ttfamily] at (3.2,\y) {tracker.example};
    \draw[msg, netred] (4.5,\y) -- node[above, font=\scriptsize\ttfamily] {Cookie: id=1678} (7.6,0);
  }
  \node[draw, fill=netorange!40, rounded corners, minimum height=0.8cm, font=\scriptsize\ttfamily] (tr) at (8.8,0) {tracker.example};
  \node[anchor=west, text width=2.6cm, text=netred] at (10.2,0) {un numero, tre siti: una cronologia};
\end{tikzpicture}
\caption{Cookie di terze parti: un solo identificatore attraverso molti siti.}
\end{figure}

Una pagina incorpora un'immagine o uno script di un altro dominio, e quel dominio imposta un cookie suo. Lo stesso cookie torna da ogni
sito che lo incorpora, insieme al sito su cui eravate. I browser oggi isolano questi cookie per sito, o li bloccano, in misura diversa.

\warning{Niente cookie, e ancora riconosciuti}{
    Cancellate ogni cookie e tornate. Il sito può ancora distinguervi da tutti gli altri? Spesso sì: user agent, dimensione dello schermo,
    font, fuso orario, lingua, il modo in cui la scheda grafica disegna un'immagine di prova: insieme sono quasi unici, un'\textbf{impronta}
    (\emph{fingerprint}). Viene inviata senza chiedere, non richiede nulla di memorizzato sulla vostra macchina, e cancellare i dati non la
    azzera. Cambiare valori in modo che nessun altro abbia è un segno a sé stante: la difesa è l'\textbf{uniformità}, ogni utente che
    riporta gli stessi valori. Tor Browser fa esattamente questo.
}

\subsection{Ogni nome, prima di ogni connessione}

Prima di una connessione c'è una query, in chiaro su UDP porta 53, con il nome che state per visitare. Attraversa la rete locale e
l'ISP, e raggiunge il resolver: per default, quello che vi ha dato il DHCP, cioè quello dell'ISP. Il resolver vede ogni nome che cercate,
dal vostro indirizzo, con l'ora, e può conservarli. DoT e DoH nascondono la query al percorso, non al resolver: deve rispondere, e quindi
legge.

\warning{DNS cifrato: chi legge ora i nomi}{
    Passate a DoH verso un grande resolver pubblico. Chi sa quali nomi cercate? Il resolver pubblico, tutti. Il vostro ISP non vede più le
    query, e un attimo dopo vede la SNI e l'indirizzo della connessione che segue. Cifrare un canale \textbf{non elimina} l'osservatore
    all'altro capo: lo \textbf{sposta}. Avete scambiato un resolver che non avevate scelto con uno che avete scelto, e ora dovete fidarvi di
    quello. Teniamolo a mente: è lo schema di ogni strumento di questo capitolo.
}

\subsection{Che cosa lascia in chiaro TLS}

\begin{itemize}
    \item \textbf{Il nome}: la SNI nel Client Hello. ECH la cifra, con una chiave che il sito pubblica nel DNS.
    \item \textbf{L'indirizzo} di destinazione, in ogni datagramma: TLS sta sopra IP e non può toccarlo.
    \item \textbf{La porta}, che di solito dice l'applicazione: 443 il Web, 853 il DNS su TLS, 993 la posta.
    \item \textbf{Dimensioni e tempi}: una pagina è una sequenza di oggetti di lunghezze note, e la sua forma sopravvive alla cifratura.
          Associare forme cifrate a pagine note è il \textbf{website fingerprinting}: funziona, meno bene, perfino attraverso Tor.
\end{itemize}

\info{Un nome nascosto ha bisogno di una folla di nomi}{
    Con ECH il nome è cifrato e resta solo l'indirizzo. Quando l'indirizzo tradisce comunque il sito? Quando il sito ha un indirizzo suo.
    Dietro una grande CDN migliaia di nomi condividono gli stessi indirizzi, e l'indirizzo dice solo quale CDN. ECH nasconde un nome solo
    \textbf{tra gli altri nomi} dietro la stessa facciata, e cambia anche l'aritmetica di un censore: bloccare un sito significa bloccarli
    tutti. Entrambe le idee, nascondersi in una folla ed essere costosi da bloccare, tornano con Tor.
}

\subsection{L'indirizzo IP}

Perché non scrivere semplicemente un indirizzo sorgente falso? Perché è anche l'\textbf{indirizzo di ritorno}: la risposta va a chi
possiede quello falsificato, e la connessione non si completa mai. Falsificare funziona solo quando non si vuole risposta (gli attacchi
per riflessione); TCP ha bisogno del SYN-ACK, e ogni richiesta si aspetta un contenuto in cambio. Per nascondere l'indirizzo di ritorno
serve qualcuno che riceva la risposta al posto vostro e ve la passi: un \textbf{proxy}.

\begin{itemize}
    \item Un indirizzo rimanda a un \textbf{contratto}: l'ISP l'ha assegnato e sa quale abbonato lo aveva in quale momento. Un tribunale può
          chiederlo.
    \item I blocchi di indirizzi sono registrati per regione: un indirizzo pubblico vi localizza, all'incirca, per chiunque.
    \item \textbf{NAT}: tutta la casa esce con un indirizzo, e siete nascosti tra i familiari, non nel mondo. Lo spazio di indirizzi
          condiviso del provider vi nasconde tra i suoi clienti, e il provider tiene la tabella.
    \item \textbf{IPv6}: un indirizzo globale per dispositivo, e nessuna folla; gli host scelgono indirizzi temporanei e continuano a
          cambiarli.
\end{itemize}

\subsection{L'indirizzo MAC e la radio}

\begin{itemize}
    \item L'indirizzo MAC non lascia mai la LAN: ogni router scarta il frame e ne costruisce uno nuovo, quindi il sito non lo vede mai. Lo
          vedono però tutti sulla vostra LAN, e chi gestisce l'access point: il bar, l'albergo, il campus.
    \item Un MAC fisso riconosce lo stesso dispositivo in ogni rete a cui si unisce: da qui gli indirizzi locali casuali che abbiamo
          incontrato. Molti sistemi tengono un indirizzo casuale per rete, così che quella rete vi riconosca quando tornate. E la richiesta
          DHCP porta spesso il nome della macchina, come l'ha impostato il proprietario.
    \item Laggiù la perdita è fisica, non un header: per essere ricevuti bisogna irradiare. La potenza cala con la distanza, quindi
          ricevitori che confrontano intensità e direzione localizzano un trasmettitore. La rete cellulare deve conoscere la vostra cella
          per consegnare una chiamata; un telefono che cerca reti Wi-Fi note si annuncia a ogni access point a portata. Nessuna cifratura
          nasconde che una trasmissione è avvenuta, da dove e quando.
\end{itemize}

\begin{table}[H]
\centering
\small
\begin{tabular}{p{2cm}p{3.3cm}p{3.6cm}p{4.2cm}}
\toprule
 & \textbf{Che cosa trapela} & \textbf{Chi lo vede} & \textbf{Che cosa lo riduce} \\
\midrule
Applicazione & login, cookie, impronta & il sito, le terze parti incorporate & niente login, cookie isolati, uniformità \\
DNS & ogni nome & il percorso, il resolver & DoT, DoH: resta il resolver \\
TLS, TCP & nome, porta, dimensioni, tempi & chiunque sul percorso & ECH, in una folla \\
Rete & entrambi gli indirizzi & router, il sito, i registri dell'ISP & un proxy: l'indirizzo di qualcun altro \\
Collegamento & MAC, nome della macchina & la LAN, l'access point & un MAC casuale per rete \\
Fisico & che si trasmette, e da dove & chiunque a portata & solo il silenzio \\
\bottomrule
\end{tabular}
\caption{Chi vede che cosa, scendendo nella pila.}
\end{table}

Ogni rimedio \textbf{sposta} un osservatore, o ne restringe la vista. Nessuno lo elimina.

% ══════════════════════════════════════════════════════════════
\section{Gli strumenti}
% ══════════════════════════════════════════════════════════════

\subsection{Un hop: il proxy anonimizzante e la VPN}

\begin{figure}[H]
\centering
\begin{tikzpicture}[every node/.style={font=\scriptsize\sffamily}]
  \node (a) at (0,0) {\icon[0.8cm]{laptop}}; \node[below=-1pt of a] {Alice};
  \node[netcloud, minimum width=1.5cm, minimum height=0.9cm] (i) at (3.4,0) {ISP};
  \node[draw, fill=netorange!30, rounded corners, minimum height=0.8cm, align=center] (p) at (7,0) {proxy\\o VPN};
  \node (s) at (10.6,0) {\icon[0.45cm]{server}}; \node[below=-1pt of s] {il sito};
  \draw[line width=3pt, netgreen!60!black] (a) -- (i); \draw[line width=3pt, netgreen!60!black] (i) -- (p);
  \node[above=7pt, text=netgreen!50!black] at (3.4,0.35) {tunnel cifrato};
  \draw[msg] (p) -- node[above] {la richiesta} (s);
  \node[text=netred, align=center] at (3.4,-1.1) {Alice e il proxy;\\dimensioni, tempi};
  \node[text=netred, align=center] at (7,-1.1) {Alice, il sito, e il contenuto\\se non c'è TLS};
  \node[text=netred, align=center] at (10.6,-1.1) {l'indirizzo\\del proxy};
\end{tikzpicture}
\caption{Il proxy anonimizzante: che cosa vede ciascuno.}
\end{figure}

Alice apre TLS verso un proxy, che fa la richiesta per lei: è il proxy istituzionale, usato al contrario. Una \textbf{VPN} lo fa per tutto
il traffico: ogni datagramma incapsulato, cifrato, dentro un altro, un \textbf{tunnel}. In entrambi i casi il proxy è un \textbf{man in
the middle che avete scelto}: vede ciò che vedeva l'ISP, e di più.

\warning{Il root server che rispondeva troppo bene}{
    Dietro una certa VPN un root server ha restituito un indirizzo che non poteva avere. Chi ha risposto? Il fornitore della VPN, a nome del
    root server: sta sul percorso di tutto ciò che inviate, e può rispondere come chiunque. Avete sostituito l'ISP con un osservatore che
    vede le stesse cose e ha anche il vostro account e il vostro pagamento. ``Nessun log'' è una promessa sul suo comportamento, non un
    meccanismo che potete verificare. Funziona contro la rete locale e l'ISP; è inutile contro il fornitore stesso. Una VPN aziendale è
    un'altra cosa: esiste per farvi entrare, non per nascondervi.
}

\subsection{Tor: un circuito di tre relay}

\begin{figure}[H]
\centering
\begin{tikzpicture}[every node/.style={font=\scriptsize\sffamily}]
  \node (a) at (0,0) {\icon[0.8cm]{laptop}}; \node[below=-1pt of a] {Alice};
  \foreach \x/\n/\k/\c in {3/guard/G/netblue, 6/middle/M/netgreen!60!black, 9/exit/E/netorange} {
    \node[draw=\c, very thick, fill=white, circle, minimum size=1cm] (\k) at (\x,0) {\n};
    \node[below=2pt of \k, text=\c] {chiave \k};
  }
  \node (s) at (12,0) {\icon[0.45cm]{server}}; \node[below=-1pt of s] {il sito};
  \draw[wired] (a) -- (G); \draw[wired] (G) -- (M); \draw[wired] (M) -- (E); \draw[wired] (E) -- (s);
  \node[text width=2.7cm, align=center] at (3,1.2) {conosce Alice e il middle};
  \node[text width=2.7cm, align=center] at (6,1.2) {conosce il guard e l'exit};
  \node[text width=2.7cm, align=center] at (9,1.2) {conosce il middle e il sito};
  \node[text width=2cm, align=center] at (12,1.2) {vede l'exit};
  \begin{scope}[yshift=-1.75cm]
    \node[draw=netblue, very thick, fill=netblue!10, minimum width=1.9cm, minimum height=0.95cm] at (1.5,0) {};
    \foreach \x in {1.5,4.5} \node[draw=netgreen!60!black, very thick, fill=netgreen!10, minimum width=1.45cm, minimum height=0.7cm] at (\x,0) {};
    \foreach \x in {1.5,4.5,7.5} \node[draw=netorange, very thick, fill=netorange!15, minimum width=1cm, minimum height=0.45cm, font=\scriptsize\ttfamily] at (\x,0) {GET};
    \node[font=\scriptsize\ttfamily] at (10.5,0) {GET};
    \node[text=black!60] at (10.5,-0.45) {in chiaro, se non c'è TLS};
    \node[text=black!70] at (6,-0.9) {ogni relay toglie lo strato della propria chiave, e inoltra il resto};
  \end{scope}
\end{tikzpicture}
\caption{Un circuito Tor: tre strati di cifratura, uno per relay.}
\end{figure}

\begin{itemize}
    \item Il client sceglie tre relay gestiti da volontari, e concorda una chiave con ciascuno di essi con \textbf{Diffie-Hellman}.
    \item Avvolge la richiesta in tre strati: ogni relay toglie il proprio e inoltra il resto. È l'incapsulamento, con una chiave per strato
          (da qui il nome \emph{onion routing}).
\end{itemize}

\dfn{Perché tre relay}{
    \begin{itemize}
      \item \textbf{Uno}: il relay conosce entrambi gli estremi; è il proxy, a cui si affida tutto.
      \item \textbf{Due}: nessuno dei due conosce entrambi gli estremi, ma l'exit conosce l'ingresso: è esattamente il nodo da sorvegliare.
      \item \textbf{Tre}: l'exit conosce solo un middle, e il middle non conosce nessuno dei due estremi.
    \end{itemize}
    Tre non ferma un guard e un exit che collaborano: confrontano i tempi.
}

\subsubsection*{Tor, visto dallo stesso sniffer}

\begin{lstlisting}[language=bash]
$ sudo tshark -i enp3s0 -Y tls      # poi si avvia tor, e: torsocks curl -sI https://www.unina.it
  3 0.0402 192.168.1.23 -> 198.51.100.23 TLSv1.3 571 Client Hello (SNI=www.k3fq7zn2tb.com)
  5 0.1296 198.51.100.23 -> 192.168.1.23 TLSv1.3 1514 Server Hello, Change Cipher Spec, App...
  9 0.2615 192.168.1.23 -> 198.51.100.23 TLSv1.3 602 Application Data
 11 0.3907 198.51.100.23 -> 192.168.1.23 TLSv1.3 602 Application Data
 14 0.5213 192.168.1.23 -> 198.51.100.23 TLSv1.3 602 Application Data
 17 0.6698 198.51.100.23 -> 192.168.1.23 TLSv1.3 1204 Application Data
\end{lstlisting}

\begin{itemize}
    \item Ogni pacchetto va a un solo indirizzo, il guard. \texttt{www.unina.it} non c'è da nessuna parte: la SNI è un nome inventato da Tor.
    \item Record di poche lunghezze sempre uguali: Tor porta i dati in unità di dimensione fissa (\emph{cell}), quindi le dimensioni dicono
          molto meno.
    \item Che cosa impara ancora l'ISP: che usate Tor, quando, e quanto. L'elenco dei relay è pubblico.
\end{itemize}

\warning{L'exit legge ciò che porta}{
    Tra Tor e il sito non c'è TLS. Che cosa vede l'exit? Ciò che vedrebbe una rete dal lato del sito: la richiesta, la risposta, qualsiasi
    password vi sia dentro, ma non chi siete, a meno che il contenuto non lo dica. Gli exit sono gestiti da volontari, e alcuni sono stati
    sorpresi a riscrivere l'HTTP in chiaro di passaggio. Con TLS sopra, l'exit vede il nome e l'indirizzo, e nient'altro: ciò che vede il
    vostro ISP senza Tor. \textbf{Tor nasconde chi; TLS nasconde che cosa.} Servono entrambi.
}

\warning{Traffico attraverso Tor, nomi attraverso l'ISP}{
    Un'applicazione manda le proprie connessioni attraverso Tor ma risolve i nomi da sola. Che cosa vede l'ISP? Ogni nome, in chiaro, poco
    prima di ogni connessione: il traffico è nascosto, l'elenco delle destinazioni no. Tor risolve i nomi all'exit: il client gli consegna
    il nome, mai un indirizzo che ha cercato. E Tor porta solo TCP: ciò che un'applicazione invia su UDP passa di lato, dal vostro
    indirizzo. Una \textbf{fuga} (\emph{leak}) è qualsiasi cosa esca per un'altra strada: Tor protegge ciò che porta, e nient'altro.
}

\subsection{Un server che nessuno può localizzare}

Finora il client si nasconde e il server è pubblico. Un \textbf{onion service} nasconde anche il server.

\begin{itemize}
    \item Il server mantiene circuiti verso alcuni relay, i suoi \textbf{punti di introduzione}, e li pubblica firmati.
    \item Client e server costruiscono ciascuno un circuito verso un \textbf{relay di rendezvous}, che li unisce: sei relay, e nessun exit.
          Nessun exit significa che nessun relay tiene mai il traffico in chiaro.
    \item Il nome \texttt{.onion} codifica la chiave pubblica del server: nessun DNS da interrogare, nessuna CA di cui fidarsi.
    \item \emph{Deep} e \emph{dark}: il deep Web è ciò che i motori di ricerca non indicizzano; il dark Web, ciò che raggiungono solo reti
          come questa.
\end{itemize}

\subsection{Contro il censore: bridge e trasporti}

\begin{itemize}
    \item Per scegliere tre relay un client deve conoscerli tutti: l'elenco dei relay è pubblico per progetto. Un censore lo copia e scarta
          ogni pacchetto verso quegli indirizzi.
    \item Un \textbf{bridge} è un relay d'ingresso che non è nell'elenco, distribuito pochi alla volta. Anche il censore può chiedere bridge
          e bloccare quelli che ottiene: la distribuzione è una gara. Un bridge nasconde \emph{dove} ci si collega, non che aspetto ha la
          connessione.
    \item L'ispezione profonda riconosce l'handshake di Tor, quindi un \textbf{pluggable transport} lo fa sembrare byte casuali. Ma quasi
          tutto il resto oggi è TLS, e byte che non somigliano a niente spiccano: è la rilevazione per anomalie degli IDS. Meglio somigliare a
          qualcosa di comune, una connessione HTTPS o una videochiamata, o nascondersi dietro una CDN troppo costosa da bloccare.
          Steganografia applicata al traffico, e una corsa agli armamenti: ogni travestimento viene studiato finché non viene riconosciuto.
\end{itemize}

\subsection{I2P: il garlic routing}

\begin{table}[H]
\centering
\begin{tabular}{p{3.2cm}p{6cm}p{6cm}}
\toprule
 & \textbf{Tor} & \textbf{I2P} \\
\midrule
Chi inoltra & relay volontari; la maggior parte degli utenti si collega soltanto & ogni partecipante inoltra per gli altri \\
Percorsi & un circuito, nei due sensi & tunnel separati a senso unico, in uscita e in entrata \\
Che cosa porta uno strato & un messaggio & più messaggi insieme: un ``aglio'' (\emph{garlic}) \\
Pensato per & raggiungere l'Internet ordinaria & servizi interni alla rete \\
\bottomrule
\end{tabular}
\caption{Tor e I2P.}
\end{table}

La stessa idea, cifratura a strati su relay scelti dal client, con scommesse diverse su chi inoltra. Tunnel separati significano che la
strada di andata e quella di ritorno attraversano relay diversi.

\subsection{Proteggere la macchina, non il percorso}

Tutto ciò che precede protegge il percorso. Un problema diverso è ciò che la macchina stessa conserva e lascia uscire.

\begin{itemize}
    \item \textbf{Tails}: un sistema avviato da una chiavetta USB, che manda tutto attraverso Tor e dimentica tutto allo spegnimento.
    \item \textbf{Qubes}: ogni attività nella propria macchina virtuale; un compartimento violato non vede gli altri.
    \item \textbf{GrapheneOS}: un Android irrobustito: permessi più severi, app tenute separate, meno da attaccare.
    \item \textbf{Whonix}: due macchine; la workstation raggiunge la rete solo attraverso un gateway che esegue Tor. È di nuovo
          l'application gateway: un malware all'interno non può far trapelare il vostro indirizzo, perché non ha una rotta.
\end{itemize}

\subsection{Aggirare l'infrastruttura: le reti mesh radio}

Ogni strumento visto finora attraversa l'infrastruttura. Dove uno Stato la blocca, la si aggira.

\begin{minipage}{0.48\linewidth}
\centering
\begin{tikzpicture}[every node/.style={font=\scriptsize\sffamily}, n/.style={draw, circle, fill=white, minimum size=0.5cm, inner sep=0pt}]
  \fill[netblue!8] (0,0) circle (1.6cm); \node[text=netblue] at (-0.9,-1.25) {raggio di A};
  \node[n] (A) at (0,0) {A}; \node[n] (B) at (1.3,0.9) {B}; \node[n] (C) at (1.2,-0.9) {C};
  \node[n] (D) at (2.6,0) {D}; \node[n] (E) at (3.9,0.9) {E}; \node[n] (F) at (3.9,-0.9) {F}; \node[n] (G) at (5.2,0) {G};
  \foreach \u/\v in {A/B, A/C, B/D, C/D, D/E, D/F, E/G, F/G} \draw[netgreen!60!black, thick] (\u) -- (\v);
  \foreach \n/\h in {B/3, D/2, E/1, G/0} \node[text=netred, above right=-2pt of \n] {\h};
  \node[text width=2.6cm, align=center] at (2.6,-1.75) {niente ISP, niente router: ogni nodo inoltra};
  \node[text=netred] at (2.6,1.45) {D legge: da A a G};
\end{tikzpicture}
\end{minipage}\hfill
\begin{minipage}{0.5\linewidth}
\begin{itemize}
    \item \textbf{LoRa}: un livello fisico, \emph{chirp} a spettro allargato nelle bande ISM sotto il GHz: chilometri a bassa potenza.
    \item \textbf{LoRaWAN}: una stella sopra di esso, con dispositivi, gateway e un server: di nuovo infrastruttura.
    \item \textbf{Meshtastic}: una mesh; ogni nodo ritrasmette ciò che sente, con un limite di hop, come il TTL.
    \item Il peer-to-peer lo faceva un livello più su: niente da bloccare, e ogni peer vedeva gli indirizzi degli altri.
\end{itemize}
\end{minipage}

\warning{Che cosa costa una mesh, e che cosa espone}{
    \begin{itemize}
      \item \textbf{Banda}: un canale stretto e un segnale debole; per Shannon, da centinaia a migliaia di bit al secondo, qualche centinaio
            di byte per messaggio.
      \item \textbf{Tempo d'aria}: le bande libere limitano quanto si trasmette, in Europa circa l'1\% del tempo nella maggior parte della
            banda.
      \item \textbf{Inondazione}: ogni messaggio raggiunge ogni nodo, quindi tutta la mesh condivide un'unica capacità, che non cresce con
            essa.
      \item \textbf{Header}: ogni relay deve leggere mittente, destinazione e conteggio degli hop: chi parla con chi è pubblico a portata
            di radio. Il contenuto è cifrato con una chiave condivisa dal canale (quella predefinita è pubblica), e i radioamatori non
            possono cifrare.
      \item E irradia: utile dove uno Stato blocca; un faro dove dà la caccia ai trasmettitori.
    \end{itemize}
}

% ══════════════════════════════════════════════════════════════
\section{Ciò che nessuno strumento risolve}
% ══════════════════════════════════════════════════════════════

\subsection{Se tre relay sono buoni, perché non dieci?}

Perché l'attacco più forte ignora il centro. Chi osserva il traffico che entra al guard e quello che esce all'exit confronta i due,
qualunque sia il numero di relay in mezzo. Ogni relay in più aggiunge latenza, e un'altra probabilità di trovarne uno lento o ostile, e
non toglie nulla a quell'attacco: tre comprano ciò che i relay possono comprare.

\begin{figure}[H]
\centering
\begin{tikzpicture}[every node/.style={font=\scriptsize\sffamily}]
  \node[anchor=east, align=right] at (-0.2,0.25) {in entrata,\\al guard};
  \node[anchor=east, align=right] at (-0.2,-1.25) {in uscita,\\all'exit};
  \draw[-{Stealth}] (0,0) -- (11,0) node[right] {tempo}; \draw[-{Stealth}] (0,-1.5) -- (11,-1.5) node[right] {tempo};
  \foreach \x/\w in {0.5/0.6, 1.4/0.3, 3.2/1.1, 5.3/0.4, 6.0/0.4, 8.1/0.9} {
    \fill[netblue!60] (\x,0) rectangle ++(\w,0.5);
    \fill[netblue!60] (\x+1.2,-1.5) rectangle ++(\w,0.5);
  }
  \draw[{Stealth}-{Stealth}, netred] (0.5,-0.25) -- node[below, text=netred] {$d$} (1.7,-0.25);
  \node[text=netred] at (5.5,-2.0) {un osservatore di entrambi i collegamenti: le stesse raffiche e gli stessi silenzi, $d$ dopo};
\end{tikzpicture}
\caption{La correlazione temporale: la cifratura cambia i byte, non quando si muovono né quanti sono.}
\end{figure}

Raffiche e silenzi attraversano il circuito intatti. Rompere la corrispondenza significa trattenere, riempire e riordinare il traffico,
come facevano i remailer: le \textbf{mix network}. Costa secondi o ore. Tor ha scelto la bassa latenza, perché il Web funzioni, e per
questo può essere correlato.

\begin{itemize}
    \item Chi può vedere entrambi gli estremi? Un ISP nazionale, quando utente e sito sono nello stesso Paese; un sistema autonomo che sta
          su entrambi i percorsi (e quali lo fanno lo decide BGP, non voi); chi gestisce molti relay, e viene scelto a volte sia come
          guard sia come exit.
    \item Da qui il \textbf{guard}: un client tiene lo stesso ingresso per mesi. Scegliere ogni volta da capo rende certo, prima o poi,
          uno ostile.
    \item Tor non ha mai affermato il contrario: fin dal primo progetto, non protegge da un osservatore di entrambi gli estremi.
\end{itemize}

\subsection{Contenuti cifrati, conversazioni visibili}

\begin{itemize}
    \item Signal e Matrix cifrano i messaggi \textbf{da un estremo all'altro}: il server conserva ciò che non può leggere. Ma li consegna,
          quindi vede chi si collega, quando, quanto spesso, quanto, e di solito a chi.
    \item Signal è costruito per conservare il meno possibile, e per nascondere perfino il mittente al proprio server. La rete aggiunge la
          sua parte: il vostro indirizzo raggiunge il server con ogni messaggio.
    \item Parlare con Alice mette Bob nella lista: gli Stati repressivi capiscono la chiusura transitiva. Per molta sorveglianza
          \textbf{chi parla con chi} basta, ed è molto più economico da conservare del contenuto (i \textbf{metadati}).
    \item \textbf{Centralizzato}: un unico operatore vede i metadati di tutti. \textbf{Federato}: come la posta, ogni utente su un server, e
          i server parlano tra loro (Matrix è federato). Il vostro server tiene i vostri metadati, finché non entrate in una stanza: ogni
          server con un membro ne tiene una copia.
\end{itemize}

\warning{Tor e una carta di credito}{
    Raggiungete un negozio attraverso Tor e pagate con la carta di credito. Quanto è stato anonimo l'acquisto? Per niente: la carta porta il
    vostro nome. Pagate in bitcoin, e il negozio riceve un indirizzo su un registro pubblico per sempre: ogni pagamento da e verso
    quell'indirizzo è lì, leggibile da chiunque, e un solo collegamento con una persona dà un nome all'intera storia. Il collegamento di
    solito si fa dove le monete sono state comprate: un exchange che ha chiesto un documento. Le criptovalute sono pseudonime, come i
    cookie: date il vostro nome una volta, e ogni pagamento passato lo acquista. Monero nasconde mittente, destinatario e importo per
    default; Zcash quando l'utente lo sceglie.
}

\subsection{L'anonimato ha bisogno di una folla}

\begin{itemize}
    \item L'anonimato è sempre tra altri, l'insieme di anonimato. Da soli, l'unico candidato siete voi.
    \item Tor ha bisogno di molti utenti, e diversi: se lo usassero solo i dissidenti, usarlo sarebbe una confessione.
    \item Un server VPN tutto vostro ha un utente: tutto ciò che esce dal suo indirizzo siete voi. Una VPN commerciale ha una folla, ma il
          fornitore ci vede attraverso.
    \item Tor Browser rende i suoi utenti identici: la risposta all'impronta, e il motivo per non personalizzarlo.
    \item Ogni utente che si aggiunge protegge gli altri: l'anonimato è un \textbf{bene collettivo}.
\end{itemize}

\subsection{VPN e Tor: in quale ordine}

\begin{itemize}
    \item \textbf{Prima la VPN, poi Tor}: nasconde Tor al vostro ISP (che vede solo una VPN); il fornitore impara che lo usate. Un bridge fa
          lo stesso senza fornitore.
    \item \textbf{Prima Tor, poi la VPN}: il sito vede sempre lo stesso indirizzo; la VPN è un exit che sa chi la paga.
    \item Nessuno dei due è più sicuro in generale. La domanda è la prima del capitolo: \textbf{da chi?}
\end{itemize}

\begin{table}[H]
\centering
\begin{tabular}{lcccc}
\toprule
\textbf{Alice non vuole che} & \textbf{HTTP} & \textbf{HTTPS} & \textbf{VPN} & \textbf{Tor} \\
\midrule
il sito conosca il suo indirizzo IP & esposto & esposto & nascosto & nascosto \\
il suo ISP sappia quale sito visita & esposto & esposto & nascosto & nascosto \\
il suo ISP legga ciò che scambia & esposto & nascosto & nascosto & nascosto \\
un'unica parte conosca entrambi gli estremi & esposto & esposto & esposto & nascosto \\
un osservatore di entrambi gli estremi li colleghi & esposto & esposto & esposto & esposto \\
\bottomrule
\end{tabular}
\caption{Dipende da chi ci si nasconde.}
\end{table}

La tabella si legge per righe: ogni riga è un avversario, e ciascuno richiede la propria risposta. Un censore non c'è: contro un censore,
bridge, travestimenti, e dove questi falliscono, una mesh.

\begin{esercizio}{Un Tor con due relay}
Alice vuole una pagina di \texttt{activist.com} attraverso due proxy che non colludono, Proxy1 e Proxy2. Ha i loro certificati, con le
chiavi pubbliche $K_1^+$ e $K_2^+$. Disegnare ogni risposta come un diagramma temporale.
\begin{enumerate}[label=(\alph*)]
    \item Un protocollo, il più semplice possibile, con cui Alice condivide una chiave di sessione $S_1$ con Proxy1.
    \item Uno con cui condivide $S_2$ con Proxy2 senza rivelare a Proxy2 il proprio indirizzo IP.
    \item Con $S_1$ e $S_2$ già stabilite e senza operazioni a chiave pubblica: una richiesta della pagina che nasconda l'indirizzo di Alice a
          Proxy2 e il sito a Proxy1, fino all'arrivo della richiesta HTTP ad \texttt{activist.com}.
    \item Per l'argomento di questo capitolo, che cosa aggiungerebbe un terzo proxy? E che cosa si perde se la chiave privata di Proxy1 viene
          rubata l'anno prossimo?
\end{enumerate}

\svolgimento
\begin{center}
\begin{tikzpicture}[yscale=0.8, every node/.style={font=\scriptsize\sffamily}]
  \timeline{al}{0}{Alice}{5.2}
  \timeline{p1}{3.6}{Proxy1}{5.2}
  \timeline{p2}{7.2}{Proxy2}{5.2}
  \timeline{si}{10.4}{activist.com}{5.2}
  \draw[msg, netblue] (0,-0.5) -- node[above] {$K_1^+(S_1)$} (3.6,-0.9);
  \draw[msg, netgreen!60!black] (0,-1.7) -- node[above] {$S_1(\text{Proxy2}, K_2^+(S_2))$} (3.6,-2.1);
  \draw[msg, netgreen!60!black] (3.6,-2.2) -- node[above] {$K_2^+(S_2)$} (7.2,-2.6);
  \draw[msg, netorange] (0,-3.4) -- node[above] {$S_1(\text{Proxy2}, S_2(m))$} (3.6,-3.8);
  \draw[msg, netorange] (3.6,-3.9) -- node[above] {$S_2(\text{activist.com}, m)$} (7.2,-4.3);
  \draw[msg, netorange] (7.2,-4.4) -- node[above] {$m$} (10.4,-4.8);
  \node[anchor=west] at (10.6,-4.8) {$m$ = il GET};
\end{tikzpicture}
\end{center}
\begin{enumerate}[label=(\alph*)]
    \item Alice sceglie $S_1$ e invia $K_1^+(S_1)$: solo Proxy1 può aprirlo.
    \item Invia $S_1(\text{Proxy2}, K_2^+(S_2))$. Proxy1 toglie il proprio strato e inoltra $K_2^+(S_2)$: Proxy2 vede l'indirizzo di Proxy1,
          non quello di Alice.
    \item $S_1(\text{Proxy2}, S_2(\text{activist.com}, \text{GET}))$. Proxy1 impara solo Proxy2, Proxy2 solo il sito; la risposta torna
          avvolta nel senso opposto, $S_2$ aggiunto da Proxy2 e $S_1$ da Proxy1, e Alice li toglie entrambi.
    \item Proxy2 conosce Proxy1: è il nodo da sorvegliare. Un terzo relay lascerebbe all'exit la conoscenza di un solo middle. E con $S_1$
          inviata sotto $K_1^+$, una chiave rubata apre ogni sessione registrata: per questo Tor usa \textbf{Diffie-Hellman} (forward
          secrecy), come TLS 1.3.
\end{enumerate}
\end{esercizio}

\Needspace{20\baselineskip}
% ══════════════════════════════════════════════════════════════
\section{Riepilogo}
% ══════════════════════════════════════════════════════════════

\riepilogo{
\begin{itemize}
    \item Nessuno è anonimo in generale, solo rispetto a qualcuno; e confidenzialità, anonimato, pseudonimato e non collegabilità sono
          quattro proprietà, non una.
    \item Ogni header è scritto per essere letto: ogni livello, fino alla radio, lascia trapelare ciò che gli serve per consegnare.
    \item Ogni strumento \textbf{sposta} l'osservatore. Una VPN lo sostituisce; Tor lo divide su tre relay che conoscono ciascuno due vicini.
    \item I censori bloccano invece di leggere: i bridge nascondono dove, i trasporti nascondono come, una mesh evita la rete.
    \item Niente ferma un osservatore di entrambi gli estremi, un server che conserva i metadati, o un pagamento a proprio nome.
    \item Abbiamo cominciato seguendo una richiesta HTTP attraverso la rete; finiamo sapendo chi può dire che era nostra.
\end{itemize}
}
